\section{Introduction}
\label{section:introduction}
The $n$-th general linear group is the group $\GL_n(\Real)$ of linear automorphisms of $\Real^n$. Thanks to the differentiability of linear maps, $\GL_n(\Real)$ comes equipped with a smooth structure that makes it into a Lie group. A Lie group is a group object in the category $\SMan$ of finitely-dimensional smooth manifolds. More concretely, a Lie group consists of a group whose underlying set is a smooth manifold and whose multiplication and inverse maps are smooth with respect to that smooth structure. Every Lie group $G$ comes with an associated Lie algebra $\g$, which is classically defined as the Lie algebra of vector fields on $G$ that are invariant under the left action $L_g(h)\=gh$ of the group. The Lie algebra $\g$ is also isomorphic to the tangent space $\T_\e G$ of $G$ at the unit.

In previous work~\cite{lanfranchi:lie-correspondence-tangent-cats}, we internalized this correspondence between Lie groups and Lie algebras in tangent categories.

\begin{content}[Tangent categories]
Introduced by Rosick\'y in~\cite{rosicky:tangent-cats} and later generalized by Cockett and Cruttwell in~\cite{cockett:tangent-cats}, tangent categories offer a minimal categorical context for differential geometry. A tangent category consists of a collection of objects, interpreted as generalized geometric spaces, a collection of morphisms, interpreted as smooth functions, and an assignment, $\T$, which sends each object $M$ to a new object $\T M$ and each morphism $f\colon M\to N$ to a new morphism $\T f\colon\T M\to\T N$.

The object $\T M$ should be regarded as the space of tangent vectors of $M$, known as the \textbf{tangent bundle} of $M$, while $\T f$ is interpreted as the smooth function which sends each tangent vector $u\in\T M$ of $M$ at $x$ to the tangent vector $\d_xf(u)$, where $\d_xf$ denotes the differential of $f$ at $x$. The assignment $\T$ is functorial, in that $\T$ preserves identities and composition of morphisms.

A tangent category also comes with a family of natural transformations
\begin{align*}
&p_M\colon\T M\to M         &&z_M\colon M\to\T M            &&s_M\colon\T_2M\to\T M\\
&l_M\colon\T M\to\T\T M     &&                              &&c_M\colon\T\T M\to\T\T M
\end{align*}
natural in $M$, respectively called the \textbf{projection}, the \textbf{zero}, and the \textbf{sum}, which equip the tangent bundle with an additive structure on the fibre, the \textbf{vertical lift}, and finally, the \textbf{canonical flip}. The projection $p_M\colon\T M\to M$ is interpreted as the map that sends each tangent vector to its base point; the zero $z_M\colon M\to\T M$ corresponds to the map that sends each point $x$ of $M$ to the zero tangent vector $0_x$ at $x$; the sum $s_M\colon\T_2M\to\T M$ sums together two tangent vectors $u$ and $v$ that share the same base point $x$, where $\T_2M$ denotes the pullback
\begin{equation*}
% https://q.uiver.app/#q=WzAsNCxbMSwwLCJcXFQgTSJdLFsxLDEsIk0iXSxbMCwxLCJcXFQgTSJdLFswLDAsIlxcVF8yTSJdLFszLDIsIlxccGlfMSIsMl0sWzIsMSwicF9NIiwyXSxbMCwxLCJwX00iXSxbMywwLCJcXHBpXzIiXSxbMywxLCIiLDEseyJzdHlsZSI6eyJuYW1lIjoiY29ybmVyIn19XV0=
\begin{tikzcd}
{\T_2M} & {\T M} \\
{\T M} & M
\arrow["{\pi_2}", from=1-1, to=1-2]
\arrow["{\pi_1}"', from=1-1, to=2-1]
\arrow["\lrcorner"{anchor=center, pos=0.125}, draw=none, from=1-1, to=2-2]
\arrow["{p_M}", from=1-2, to=2-2]
\arrow["{p_M}"', from=2-1, to=2-2]
\end{tikzcd}
\end{equation*}
and is a tangent pullback, in that it is preserved by all iterates $\T^n\=\T\T{\dots}\T$ of the tangent bundle functor. It is sometimes asked that $\T M$ be equipped with a \textbf{negation}, $n_M\colon\T M\to\T M$, also natural in $M$, which sends each tangent vector $u$ to $-u$. In that case, we say that the tangent category \textbf{has negatives}.

The vertical lift $l_M\colon\T M\to\T\T M$ encodes \emph{local linearity} of the tangent bundle by turning the following diagram
\begin{equation*}
% https://q.uiver.app/#q=WzAsNCxbMSwwLCJcXFRcXFQgTSJdLFsxLDEsIlxcVCBNIl0sWzAsMSwiTSJdLFswLDAsIlxcVF8yTSJdLFszLDIsIlxccGlfMXBfTSIsMl0sWzIsMSwiel9NIiwyXSxbMCwxLCJcXFQgcF9NIl0sWzMsMCwiXFx4aV9NIl1d
\begin{tikzcd}
{\T_2M} & {\T\T M} \\
M & {\T M}
\arrow["{\xi_M}", from=1-1, to=1-2]
\arrow["{\pi_1p_M}"', from=1-1, to=2-1]
\arrow["{\T p_M}", from=1-2, to=2-2]
\arrow["{z_M}"', from=2-1, to=2-2]
\end{tikzcd}
\end{equation*}
into a pullback, where:
\begin{align*}
&\xi_M\=(z_{\T M}\times l_M)\T s_M\colon\T_2M\to\T\T M
\end{align*}
This universal condition makes the fibres $\T_xM\=p^{-1}(x)$ of the tangent bundle $\T M$ of $M$ into \textit{linear spaces}, in that the tangent bundle $\T\T_xM$ of $\T_xM$ is isomorphic to the Cartesian product $\T_xM\times\T_xM$. This linearity condition is precisely the condition enjoyed by Euclidean spaces, that are $\Real^n$, for which $\T\Real^n\cong\Real^n\times\Real^n$. This expresses the idea that an object in a tangent category is a \emph{locally linear} space and a morphism is a \emph{locally linear} function.

Lastly, the canonical flip, $c_M\colon\T\T M\to\T\T M$, is an isomorphism of the double tangent bundle that encodes the symmetry of partial derivatives, in that $\partial_i\partial_jf=\partial_j\partial_if$.

The archetypal example is the category of finitely-dimensional smooth manifolds with the usual notion of the tangent bundle functor. However, tangent categories capture geometric models well beyond smooth manifolds. In~\cite{cruttwell:algebraic-geometry}, the category of affine schemes was shown to carry a tangent structure which captures some key aspects of algebraic geometry. This construction was later extended in~\cite{ikonicoff:operadic-algebras-tangent-cats} to algebras over an arbitrary algebraic operad. Convenient manifolds and $C^\infty$-rings also form tangent categories as well as the subcategory of microlinear objects of a model of Synthetic Differential Geometry (SDG) and every Cartesian Differential Category (CDC)~\cite{cockett:tangent-cats}. More recently, in~\cite{blohmann:elastic-diffeological-spaces}, Blohmann showed that \emph{elastic} diffeological spaces also form a tangent category.
\end{content}

\begin{content}[Lie groups in tangent categories]
A \textbf{Lie group} (\cite[Definition~2.13]{lanfranchi:lie-correspondence-tangent-cats}) in a Cartesian tangent category, that is, a tangent category $(\X,\TT)$ with terminal object $\*$ and binary products $A\times B$ preserved by the tangent bundle functor, is a group object $G$ with unit map $\e\colon\*\to G$, multiplication $\m\colon G\times G\to G$, and inverse map $\i\colon G\to G$, of $(\X,\TT)$, which admits tangent space at the unit, in that the pullback
\begin{equation*}
% https://q.uiver.app/#q=WzAsNCxbMSwwLCJcXFQgRyJdLFsxLDEsIkciXSxbMCwwLCJcXFRfXFxlIEciXSxbMCwxLCJcXCoiXSxbMywxLCJcXGUiLDJdLFswLDEsInBfRyJdLFsyLDNdLFsyLDAsIlxcaW90YV9cXGUiXSxbMiwxLCIiLDEseyJzdHlsZSI6eyJuYW1lIjoiY29ybmVyIn19XV0=
\begin{tikzcd}
{\T_\e G} & {\T G} \\
{\*} & G
\arrow["{\iota_\e}", from=1-1, to=1-2]
\arrow[from=1-1, to=2-1]
\arrow["\lrcorner"{anchor=center, pos=0.125}, draw=none, from=1-1, to=2-2]
\arrow["{p_G}", from=1-2, to=2-2]
\arrow["\e"', from=2-1, to=2-2]
\end{tikzcd}
\end{equation*}
of the tangent bundle $p_G\colon\T G\to G$ along the unit map $\e$ exists and is preserved by all iterates $\T^n$ of the tangent bundle functor.
\end{content}

\begin{content}[Lie functors]
Every group object $G$ in a Cartesian tangent category whose tangent bundle $p_G$ has negation $n_G\colon\T G\to\T G$, comes with a functor
\begin{align*}
&\g(-)\colon\X\to\LieAlg^\op
\end{align*}
called the \textbf{Lie functor} of $G$, that sends every object $A$ of $\X$ to the Lie algebra $\g(A)$ of \textbf{$A$-Parametrized Left-Invariant Vector Fields} (\textbf{$A$-\PLIVFs}\!) of $G$ (\cite[Definition~3.18]{lanfranchi:lie-correspondence-tangent-cats}), which are maps $X\colon G\times A\to\T G$ that satisfy the following two conditions:
\begin{equation*}
% https://q.uiver.app/#q=WzAsMyxbMCwwLCJHXFx0aW1lcyBBIl0sWzEsMCwiXFxUIEciXSxbMSwxLCJHIl0sWzAsMSwiWCJdLFsxLDIsInBfRyJdLFswLDIsIlxccGlfMSIsMl1d
\begin{tikzcd}
{G\times A} & {\T G} \\
& G
\arrow["X", from=1-1, to=1-2]
\arrow["{\pi_1}"', from=1-1, to=2-2]
\arrow["{p_G}", from=1-2, to=2-2]
\end{tikzcd}\qquad
% https://q.uiver.app/#q=WzAsNCxbMSwwLCJHXFx0aW1lcyBBIl0sWzEsMSwiXFxUIEciXSxbMCwwLCJHXFx0aW1lcyBHXFx0aW1lcyBBIl0sWzAsMSwiXFxUKEdcXHRpbWVzIEcpIl0sWzAsMSwiWCJdLFsyLDAsIlxcbVxcdGltZXNcXGlkX0EiXSxbMiwzLCJ6X0dcXHRpbWVzIFgiLDJdLFszLDEsIlxcVFxcbSIsMl1d
\begin{tikzcd}
{G\times G\times A} & {G\times A} \\
{\T(G\times G)} & {\T G}
\arrow["{\m\times\id_A}", from=1-1, to=1-2]
\arrow["{z_G\times X}"', from=1-1, to=2-1]
\arrow["X", from=1-2, to=2-2]
\arrow["{\T\m}"', from=2-1, to=2-2]
\end{tikzcd}
\end{equation*}
The Lie functor $\g(-)$ of a Lie group is always representable (\cite[Theorem~4.20]{lanfranchi:lie-correspondence-tangent-cats}), and it is represented by the tangent space $\T_\e G$ of $G$. Thus, via the Yoneda lemma, the Lie algebra structure of each $\g(A)$ induces an \emph{internal} Lie algebra structure on $\T_\e G$ (\cite[Theorem~5.4]{lanfranchi:lie-correspondence-tangent-cats}).
\end{content}

\begin{content}[Differential Lie algebras]
The resulting Lie algebra object, denoted by $\g\in(\X,\TT)$, becomes a \textbf{differential Lie algebra} (\cite[Definition~5.1 and Theorem~5.12]{lanfranchi:lie-correspondence-tangent-cats}), in that the underlying object of $\g$ is a differential object, the additive structure of the Lie algebra coincides with that of the differential object $\g$, and the Lie bracket $[,]\colon\g\times\g\to\g$ is bilinear, that is:
\begin{equation*}
% https://q.uiver.app/#q=WzAsNixbMCwxLCJcXGdcXHRpbWVzXFxnIl0sWzEsMSwiXFxnIl0sWzAsMCwiXFxUKFxcZ1xcdGltZXNcXGcpIl0sWzEsMCwiXFxUKFxcZykiXSxbMiwxLCJcXGdcXHRpbWVzXFxnIl0sWzIsMCwiXFxUKFxcZ1xcdGltZXNcXGcpIl0sWzIsMywiXFxUWyxdIl0sWzAsMSwiWyxdIiwyXSxbMCwyLCJ6X1xcZ1xcdGltZXNcXGxhbWJkYV9cXGciXSxbMSwzLCJcXGxhbWJkYV9cXGciLDJdLFs1LDMsIlxcVFssXSIsMl0sWzQsMSwiWyxdIl0sWzQsNSwiXFxsYW1iZGFfXFxnXFx0aW1lcyB6X1xcZyIsMl1d
\begin{tikzcd}
{\T(\g\times\g)} & {\T(\g)} & {\T(\g\times\g)} \\
{\g\times\g} & \g & {\g\times\g}
\arrow["{\T[,]}", from=1-1, to=1-2]
\arrow["{\T[,]}"', from=1-3, to=1-2]
\arrow["{z_\g\times\lambda_\g}", from=2-1, to=1-1]
\arrow["{[,]}"', from=2-1, to=2-2]
\arrow["{\lambda_\g}"', from=2-2, to=1-2]
\arrow["{\lambda_\g\times z_\g}"', from=2-3, to=1-3]
\arrow["{[,]}", from=2-3, to=2-2]
\end{tikzcd}
\end{equation*}
where $\lambda_\g\colon\g\to\T\g$ is the vertical lift of the differential object $\g$.

Recall that a \textbf{differential object} (\cite[Definition~3.1]{cockett:differential-bundles}) in a Cartesian tangent category consists of a commutative monoid $(\A,0,+)$ equipped with a \textbf{vertical lift} $\lambda_\A\colon\A\to\T\A$, compatible with the additive structure, that makes the following morphism
\begin{align*}
&\hat\xi_\A\colon\A\times\A\xrightarrow{z_\A\times\lambda_\A}\T(\A\times\A)\xrightarrow{\T+}\T\A
\end{align*}
into an isomorphism. Equivalently, a differential object consists of a commutative monoid $(\A,0,+)$ equipped with a \textbf{differential projection}, which is a map $\hat p_\A\colon\T\A\to\A$ that is compatible with the additive structure of $\A$ and of its tangent bundle, and that turns the following diagram
\begin{equation*}
% https://q.uiver.app/#q=WzAsMyxbMSwwLCJcXFRcXEEiXSxbMCwwLCJcXEEiXSxbMiwwLCJcXEEiXSxbMCwxLCJwX1xcQSIsMl0sWzAsMiwiXFxoYXQgcF9cXEEiXV0=
\begin{tikzcd}
\A & {\T\A} & \A
\arrow["{p_\A}"', from=1-2, to=1-1]
\arrow["{\hat p_\A}", from=1-2, to=1-3]
\end{tikzcd}
\end{equation*}
into a product diagram.
\end{content}

\begin{content}[Goal of the paper]
In this paper, we introduce and study general linear group objects and their Lie algebras in tangent categories. We believe this work is not only important in itself since it introduces an important class of group objects in the general context of tangent categories, but it is also preliminary to developing a rich theory of geometric constructions in the tangent-categorical world. In particular, general linear groups are an important ingredient in the definition of the \emph{frame} bundle of a vector bundle $\q\colon E\to M$, which is defined as a $\GL_n(\Real)$-principal bundle, where $n$ is the dimension of the typical fibre of $\q$.
\end{content}

\begin{content}[Structure of the paper]
The paper is divided into three major parts. In the first part, we take a deep dive into multilinearity in the context of tangent categories and, in particular, in the notion of internal multilinear Homs (Definition~\ref{definition:linearly-closed-tangent-category}), which provide internal Homs for multilinear maps between differential objects. While the notions of multilinearity are not new in the tangent-categorical literature (see~\cite{cruttwell:sector-forms-in-CDCs} and~\cite{lemay:jacobian}), so far, it has only been considered in the context of Cartesian differential categories. In this part, we will also introduce \textbf{Linear Endomorphism Classifiers} (Definition~\ref{definition:linear-endomorphism-classifier}), which are those internal multilinear Homs of type $\lnr[\A,\A]$. We will also study their algebraic properties and show that their maximal differential subring $\lnr_\A^-$ (Definition~\ref{definition:maximal-differential-subring}) exhibits the structure of a differential Lie algebra.

In Section~\ref{section:general-linear-groups}, we make use of the knowledge we gained in the first section about internal multilinear Homs and introduce general linear group objects, as classifiers of a special class of maps that we call \textbf{Kleisli parametrized linear automorphisms} (Definition~\ref{definition:kleisli-PLAs}), which are maps $f\colon\A\times M\to\T^n\A$ that lift to the Kleisli tangent category $\KL[\T_M^n,c_M^n]$ as linear automorphisms $f\colon\A\nto\A$ (Lemma~\ref{lemma:kleisli-PLI}). After introducing general linear groups $\GL_\A$ of each differential object $\A$, we show how to construct them directly from the linear endomorphism classifier $\lnr_\A$ via a tangent triple equalizer (Theorem~\ref{theorem:triple-equalizer}).

In the third and last part, we study the differential Lie algebras of the general linear groups. We begin by constructing a natural Lie algebra homomorphism $\chi_M\colon\gl_\A(M)\to\lnr_\A^-(M)$ between the Lie functor of $\GL_\A$ and the functor $\lnr_\A^-(-)$, which is represented by the Lie algebra $\lnr_\A^-$ of the maximal subring of the linear endomorphism classifier (Proposition~\ref{proposition:the-map-chi}). Using this construction and the characterization of Theorem~\ref{theorem:triple-equalizer}, we prove the main result of the paper: that the differential Lie algebra $\gl_\A$ of the general linear group $\GL_\A$ of a differential object $\A$ is isomorphic to the differential Lie algebra $\lnr_\A^-$ (Theorem~\ref{theorem:main-theorem}). Finally, when $\A$ admits negatives, this theorem recaptures the classic result of differential and algebraic geometry, that $\gl_\A\cong\lnr_\A$ (Corollary~\ref{corollary:main}).
\end{content}

\begin{conventions}
We shall denote composition diagrammatically; that is, $fg$ is the composition of $f$ followed by $g$. When considering functors or concrete functions, we switch to the functional notation, so for example, we shall denote by $GFM$ or $GFf$ the \emph{application} of the functor $G$ to the application of the functor $F$ to the object $M$ and to the morphism $f$.

We shall denote generic differential objects as $\A,\B,\C$ and we shall denote the zero, sum, vertical lift, and differential projection of $\A$ by $0_\A$, $+_\A$, $\lambda_\A$, and $\hat p_\A$. We shall omit the subscript $\A$ when it will be clear from the context.
\end{conventions}



%__________________________________________________________________________
%__________________________________________________________________________